%%%%%%%%%%%%%%%%%%%%%%%%%%%%%%%%%%%%%%%%%%%%%%%%%%%%%%%%%%%%%%%%%%%%%
%                                                                   %
%                           TESI in LaTex                           %
%                                                                   %
%%%%%%%%%%%%%%%%%%%%%%%%%%%%%%%%%%%%%%%%%%%%%%%%%%%%%%%%%%%%%%%%%%%%%



%%%%%%%%%%%%%%%%%%%%%%%%%%%%%%%%%%%%%%%%%%%%%%%%%%%%%%%%%%%%%%%%%%%%%
%                             Preambolo                             %
%%%%%%%%%%%%%%%%%%%%%%%%%%%%%%%%%%%%%%%%%%%%%%%%%%%%%%%%%%%%%%%%%%%%%


% Documento pronto da stampare
\documentclass[a4paper,12pt,twoside]{report}
% Documento di brutta (al poste delle immagine una X)
%\documentclass[draft,a4paper,12pt,oneside]{report}


% Imposto la lingua italiana
\usepackage[italian]{babel}


% Settaggio dei margini
\topmargin 0cm
\headsep 1cm
\headheight 0.6cm
\textwidth 14.6cm
\textheight 21.8cm
\evensidemargin 1cm
\oddsidemargin 1cm


% Pacchetti da includere per...

% simboli matematici
\usepackage{amsmath}

% simboli vari
\usepackage{textcomp}
\usepackage[mathscr]{euscript}

% figure
\usepackage{epsfig}
\usepackage{graphicx}
\usepackage{xcolor}
\usepackage{subfigure}

% intestazioni
\usepackage{fancyhdr}

% encoding dei caratteri
\usepackage[ansinew]{inputenc}
\usepackage[OT1]{fontenc}

% ridefinire le didascalie
\usepackage[footnotesize,hang,bf]{caption}

% Stile della pagina
\fancyhead[L]{\chaptername \thechapter}
\fancyhead[LO]{\thesection} \fancyhead[RO]{\sectionmark}
\lhead{\slshape Ch. \thechapter}


% Configurazioni varie

% nonumchapter serve per introduzioni, prefazioni, ecc.
% titola come un capitolo ma senza numero
\def\nonumchapter#1{
    \chapter*{#1}
    \addcontentsline{toc}{chapter}{#1}
}

% estensioni e cartella delle immagini (non necessario...)
%\DeclareGraphicsExtensions{gif,eps}
%\graphicspath{{./eps/}}

% crea una nuova lunghezza e gli assegna un valore
\newlength{\defbaselineskip}
\setlength{\defbaselineskip}{\baselineskip}

% comando per settare  la variabile \baselineskip come multiplo di \defbaselineskip
\newcommand{\setlinespacing}[1]%
           {\setlength{\baselineskip}{#1 \defbaselineskip}}

% permette l'inserimento "rapido" del virgolettato (ad es. per citazioni)
\newcommand{\virgolette}[1]{``#1''}

% definizione dello stile delle didascalie delle immagini
\renewcommand{\captionfont}{%
    \normalfont \sffamily \slshape \footnotesize%
}

% questo ambiente serve per i log e codice sorgente:
% riduce il carattere la dimensione dei font e imposta l'ambiente
% come verbatim per evitare l'interpretazione di caratteri speciali
\newenvironment{logs}{\begin{flushleft}\begin{ttfamily}\scriptsize\setlinespacing{1}}{\end{ttfamily}\end{flushleft}}



%%%%%%%%%%%%%%%%%%%%%%%%%%%%%%%%%%%%%%%%%%%%%%%%%%%%%%%%%%%%%%%%%%%%%
%                         Inizio Documento                          %
%%%%%%%%%%%%%%%%%%%%%%%%%%%%%%%%%%%%%%%%%%%%%%%%%%%%%%%%%%%%%%%%%%%%%


\begin{document}


% Inizio del Frontespizio

\thispagestyle{empty}
\enlargethispage{60mm}
\begin{center}
\Large{\textsc{Politecnico di Milano}}\\
%\vspace{5mm}
\large{V Facolt\`a di Ingegneria}\\
%\vspace{5mm}
\large{Corso di laurea in Ingegneria delle Telecomunicazioni}\\
\large{Dipartimento di Elettronica e Informazione}\\
\vspace{7mm}
\begin{figure}[h]
\begin{center}
\includegraphics[scale=0.15]{logo-polimi.eps}
\end{center}
\end{figure}
\vspace{2mm}

% titolo della tesi
\begin{LARGE}
TITOLO DELLA TESI
\end{LARGE}
\vspace{25mm}

% relatore
\begin{flushleft}
\begin{tabular}{l l }
Relatore:    & Prof. Nome COGNOME\\
Correlatore: & Ing. Nome COGNOME\\
\end{tabular}
\end{flushleft}
\vspace{25mm}

% autore/autori
\begin{flushright}
\begin{tabular}{l l }
Tesi di Laurea di: & \\
Nome COGNOME & Matr. xxxxxx
\end{tabular}
\end{flushright}
\vspace{43mm}
{\large{\bf Anno Accademico 200x-200x}}
\end{center}

% Fine del Frontespizio


\clearpage
\newpage


% Settaggio interlinea
\setlinespacing{1.5}


% Inizio Numerazione Romana
\pagenumbering{Roman}


% Indice - Per abilitarlo decommentare le tre linee seguenti
%\tableofcontents
%\newpage
%\clearpage


% Elenco delle Figure
% La linea seguente inserisce la lista delle figure come voce nell'indice
%\addcontentsline{toc}{section}{\listfigurename}

% La linee seguenti inseriscono la lista delle figure nel corpo
%\listoffigures
%\clearpage


% Impostazioni della pagina
\pagestyle{fancy}
\headsep=40pt
\lhead{}
\rhead{\slshape \leftmark}
\cfoot{\thepage}


% Sommario (1 pagina)
% [Includo il file abstract.tex]
\include{abstract}


% Inizio Numerazione Araba
\pagenumbering{arabic}


% Vari Capitoli
% [Includo i vari file chapterN.tex]

%\include{introduzione}      % Introduzione

%\include{chapter2}          % Capitolo 2

%\include{chapter3}          % Capitolo 3

%\include{chapter4}          % Capitolo 4

%\include{chapter5}          % Capitolo 5

%\include{chapter6}          % Capitolo 6

%\include{chapter7}          % Capitolo 7

% Bibliografia
% [Includo il file biblio.tex.... se non usi bibtex]
%\input{biblio}


% Fine del documento

\end{document}
